% CHAPTER 7: DISCUSSION
% ======================================================================
\section{Discussion}

\subsection{Challenges and Resolutions}

Frontend-Backend API Contract Alignment. During early integration, there was plenty of mismatching response structure between Vue 3 frontend API and SpringBoot backend, which caused numerous runtime errors and slowed down development iteration speed. The resolution was standardizing on a Result\textless T\textgreater generic wrapper (code, msg, data) for all microservice responses. Frontend Axios interceptors were configured to automatically unwrap the data field, so that component code only needs to handle domain data without manually parsing response layers. Additionally, Mock detect fallback is implemented on the backend so that frontend development could proceed independently through mock data without complete backend.

LLM API Integration for AI Intent Routing. The Tool Agent relies on the DeepSeek LLM for natural language intent classification, requiring prompt engineering to produce consistent and parseable output. The team has developped several iterations of prompt design. Early versions produced verbose natural language responses that couldn't be reliably parsed, causing the intent routing switch statement to fail or select the wrong handler. A structured output format with explicit categories and parameter schemas was eventually adopted, combined with a mock-detect fallback that returns one of three hardcoded routes when the API is unreachable. Since the iFlytek MaaS platform is compatible with the OpenAI interface, the existing OpenAI library can be used directly.

Copilot Intent Classification. The Copilot adds a second layer of intent classification on top of the Tool Agent's own routing. Getting the classifier to reliably distinguish between CODE, TOOL, and RAG intents required careful prompt design. The key insight was to keep the classification prompt minimal, the team decided six category definitions, temperature 0.1, max\_tokens 10. These settings would ensure LLM returns a single word rather than a justification. Early attempts that asked the model to explain its reasoning produced verbose output that was both slower and harder to parse. The three-tier strategy (blocklist $\rightarrow$ keyword $\rightarrow$ LLM) also emerged from practical iteration: the blocklist and keyword tiers were added after observing that simple greetings and obvious jailbreak attempts were consuming LLM tokens unnecessarily. The CLARIFY intent was a later addition, added when testing showed that users frequently submitted underspecified requests (``Book a room'' without a time, ``Plan a route'' without a destination) and expected the system to ask for the missing information rather than fail silently.

RAG Permission and Vector-Space Consistency. The RAG Agent introduced various challenges from the Code and Tool agents: semantic relevance cannot be treated as authorization. In our early retrieval tests, Milvus could return highly similar data from documents which user do not have privilege to access. The team applied post-retrieval in application layer to solve this, this machanism will filter documents depend on white lists. So only authenticated chunks will be returned to LLM. Another problem is that post-retrieval would filter too much documents, so only little chunks can be returned. The team used an over-fetching strategy with a 4 x multiplier on topK, trading retrieval redundancy for result quality. This approach retrieves substantially more candidate documents from the Milvus vector store than the final required number. Even if permission filtering eliminates three-quarters of the candidates, the remaining quarter still provides enough relevant chunks to LLM.

Amap API Multi-Level Geocoding Fallback. The team found than geocoding sometimes fails for an informal address description. A three-level rollback strategy is adopted: first direct geocoding via /v3/geocode/geo, then POI text search via /v3/place/text for informal place names (e.g. "HKU", "天安门广场"), and finally a city-prefix retry. Another problem is sometimes Amap driving APi may return a null route or abnormal response. The team developped a new system. Through concatenating polyline segments from individual navigation steps, the map always has renderable coordinates.

Five-Layer SQL Validation Design. Designing a whitelist verification system requires striking a balance between security and availability. An overly strict validator will prevent legitimate analysis queries; an overly permissive one will expose the database to risk. After several rounds testing, the current five-layer architecture was found to be a good balance: operation type restriction (Layer 1) and forbidden keyword blacklist (Layer 2) provide broad security protection at the SQL system level; table and column whitelist (Layer 3 and 4) provide application-level data-scoping control; and query complexity limits (Layer 5) provide fine-grained resource governance and prevent performance-affecting Cartesian product joins. Each layer can be configured separately through the application.yml file, so different environments can adjust strictness according to their actual needs.

Model Serving Boundary. The completed implementation uses a Python HTTP inference bridge for Code Agent generation and a separate Python embedding worker for the RAG Agent. This boundary was chosen for two reasons: first, model runtime support is not equally smooth across all architectures; second, a separate Python process makes it easier to upgrade or replace models without changing the Java service or Spring context. The cost is additional network latency for HTTP calls, but the design keeps model dependencies isolated from the enterprise microservices.

\subsection{Comparison with Existing Solutions}

Compared with cloud-based LLM API services, our platform minimizes the risk of external data exposure through keeping relational storage, vector storgae, document indices, authentication and SQL execution in our own docker. The team also minimizes the usage of external API call, which is only used in intent classification, SQL and RAG generating, cloud embedding and Amap gecoding. And as for RAG, post-retrieval filtering is applied to ensure that only documents with the correct clearance level are returned to the LLM, preventing accidental leakage of sensitive information. For SQL, the database schema will be transmitted to LLM for prompt constructing. Due to the uncontrollable limitation that the team do not own high-level GPUs to train our own models, we can only use the pre-trained models as a substitute. But in real banking scenarios, the team can use the pre-trained models as a base and fine-tune them with bank-specific data to improve performance.

Compared with general multi-agent frameworks such as AutoGen and MetaGPT, our platform abandons the flexibility of dynamically allocating agent roles and temporary conversations. Instead, it adopts three long-running agents with fixed capabilities. Althpugh this design sacrifices flexibility, it provides a stable behaviour boundaries, strict access control and pridictale resource consumption, which is vital in banking scenarios.

Compared with independent RAG solutions, our integrated architecture enables unified access to three complementary AI capabilities through a single platform. The RAG component is also integrated with the same JWT, RBAC, document approval, task logging, and WebSocket progress mechanisms as the rest of the system, rather than operating as an isolated chatbot. The trade-off is increased backend complexity; each agent is a separate microservice with its own dependencies, requiring coordinated deployment through Docker Compose and a Gateway for unified entry.

\subsection{Project Limitations}

This platform still has several limitations. The Code agent's Python HTTP server approach adds network latency and requires a separate Python process which will increase memory consumption. The RAG agent is implemented end to end, but retrieval precision remains tied to document quality, chunk granularity, and the selected embedding profile, so administrators must keep the knowledge base curated. The Tool Agent depends on external APIs (DeepSeek, Amap), and the unavailability of these services would degrade functionality to the mock fallback mode. The Task Center uses in-memory queuing via LinkedBlockingQueue rather than a message broker; while adequate for the current scale, this means queued tasks are lost on service restart. The delivered deployment is single-instance and can handle approximately 50 concurrent users; multi-node Kubernetes orchestration is outside the current deployment scope. RBAC is intentionally coarse-grained, so fine-grained custom permission design is left to the adopting institution. The Copilot's intent classifier is prompt-dependent --- a significant change to the underlying LLM's behavior could require prompt adjustments, and the keyword fallback (Tier 3b) has lower accuracy than the LLM path.

\newpage

% ======================================================================
